\documentclass[a4paper,twocolumn,landscape]{article}
\usepackage{geometry}
\usepackage{concmath}
\usepackage[OT1]{fontenc}
\usepackage{graphicx}
\usepackage{amsfonts}
\usepackage{amsmath}
\usepackage{calrsfs}
\usepackage{color}
\usepackage[dvipsnames]{xcolor}
\usepackage[many]{tcolorbox}
\usepackage{setspace}              
\usepackage{multicol}  
\geometry{margin=0.5in}
\definecolor{green}{RGB}{52, 235, 168}
\definecolor{violet}{RGB}{165, 83, 237}
\definecolor{lemon}{RGB}{170, 237, 83}
\definecolor{red}{RGB}{240, 125, 84}
\definecolor{blue}{RGB}{105, 255, 242}

\begin{document}

\noindent \colorbox{green}{Ex1}\\
a/ \(f_e = \frac{1}{T_e}\) où \(f_e = 44.1kHz\)\\
donc \(T_e = 2,3 \times 10^{-5}\)
donc il y a 195 804 000 échantillons
donc 783,2 Mo\\

\noindent
b/ \(2f_{max} = 40 kHz\)\\
\(2f_{max} < f_e\) OUI\\
Avec le bruit parasite\\
\(2f_p = 84.1 kHz\)
\(2f_p > f_e\) NON\\

\noindent 
d/ Passe-bas, filtre anti-repliement \\
\(f_c = \frac{f_e}{2}\)\\

\noindent
e/ \(\frac{f_e}{2}\), 20kHz, car 20kHz \(< \frac{f_e}{2}\)\\

\noindent \colorbox{green}{Ex2}\\
1/
\(T_e = \frac{1}{f_e}\)\\
\(T_e = 50 \mu s\)\\
\(\Delta f = \frac{f_e}{N}\)\\
\(\Delta f = 20Hz\)\\
\(f_e > 2f_{max}\)\\
\(\Rightarrow f_{max} < \frac{f_e}{2}\)\\
\(f_{max} < 10kHz\)\\

\(\Delta f \searrow \Rightarrow \begin{cases}
    N \nearrow\\
    f_e \searrow
\end{cases}\\
f_{max} \nearrow \Rightarrow f_e \nearrow \)\\

\noindent 
2/ \(NT_e = t_{obs}\)\\
\(\frac{f_e}{N}<20Hz\)\\
donc \fbox{\(t_{obs} > \frac{1}{20}s\)}\\

\noindent
3/ \(f_{max} = 9f_1\)\\
\(f_{max} = 900 Hz\)\\
\fbox{\(f_e > 2f_{max}\)}
Je propose \(f_e = 2000Hz\)
il faut \(\Delta f < 2000Hz\)\\
donc \(\frac{f_e}{N}<200Hz\)\\
donc \(\frac{f_e}{200}< N\)\\
donc \fbox{\(N>10\)}\\
Je propose N = 1000\\

\noindent
\colorbox{green}{Ex3}\\
\(T_{acq} = 89s\)\\
\fbox{\(f_e = \frac{N}{T_{acq}}\)}\\
spectre 1 : \(f_e = 3.3 Hz, \frac{f_e}{2} = 1.65Hz\)\\
spectre 2 : \(f_e = 11.5Hz, \frac{f_e}{2} = 5.525Hz\)\\
spectre 3 : \(f_e = 33.3Hz, \frac{f_e}{2} = 16.6Hz\)\\
T = 0.5\\
\(f_1 = 2Hz > \frac{f_e}{2}\\
f_2 = 4Hz  \\
f_3 = 6Hz > \frac{f_e}{2}\\
f_4 = 8Hz\)\\

\noindent
\colorbox{green}{Ex4}\\
a/ \(V_A = 0 = V_M\)\\
donc quelque soit la valeur des bits 
\(I_0 = I_0'\)\\
donc \(I_1 = I_1'\)
donc \(I_k = I_k' \Rightarrow V_k\) constante\\
b/ \(V_0 = 2RI_0\)\\
\(V_1 - V_0 = RI_1 = 2RI_0 = V_0\)\\
\(V_n = 2V_0\)\\
c/ Par récurrence : \begin{equation*}
    V_{k+1} = 2V_k
\end{equation*}\\
d/  \begin{align*}
    I &= \sum_{k=0}^{n-1} a_k I_k\\
    &= \sum a_k V_k R\\
    &= \sum a_k v_0 2^k R\\
    &\Rightarrow R v_0 \sum a_k 2^k = I_0
\end{align*}
e/ Le temps de commutation des interrupteurs peut aller 
jusqu'à une dizaine de ns\\
f/ filtre passe bas caractérisé par tau
fréquence de coupure equiv temps de réponse 
\(T >> \tau >> T_e\)

\noindent \colorbox{green}{Ex5}\\
1/ \begin{align*}
    f_m = \frac{n}{NT_e} \pm \frac{1}{NT_e}\\
    \Delta f &= \frac{1}{NT_e}\\
    \Delta f \leq \frac{1}{100}f_m\\
    N \geq 250
\end{align*}
Donc \(T_e = 1ms\), \(T_{acqm} = 250ms\)\\
2/ \begin{align*}
    n_1 &= 184 \pm 1\\
    n_2 &= 368 \pm 1\\
    &= 2n_1\\
    n_3 &= 553 \pm 1
\end{align*}
pas compatible\\
3/ Dans le spectre sont présentes \(f_1, f_1 - f_2, f_1 + f_2\)
\(f_1 + f_2\) manque \(> \frac{f_e}{2}\)\\
\(f_1 = F'\)\\
\(f_2 = F'-F = 164Hz\)\\
\(f_1 + f_2 > 500Hz\)

\end{document}